% GENERATED FILE. Edit canonical lessons, then run npm run generate:books.
% canonical-source-hash: fnv1a64:36e381d3d4c6e862
% canonical-lessons: JA-C235-nageru, JA-C235-keru, JA-C235-yakyuu, JA-C235-juudou, JA-C235-yuushou, JA-R235-review-230-235

\chapter{Baseball and Judo}
\label{ch:ja-baseball-and-judo}
\hlchaptermodality{235}

\begin{chapteropening}
\textbf{By the end of this chapter:} \emph{I can talk about ball games and judo I played.}

Complete the last lesson of chapter 235: I can talk about ball games and judo I played.
\end{chapteropening}

\section[\texorpdfstring{\ja{なげる} (nageru)}{なげる (nageru)}]{\ja{なげる} (nageru) — to throw}

\label{lesson:JA-C235-nageru}



\begin{quote}
Before the new one: say the Japanese for a player, then the Japanese for cheering.
\end{quote}

\subsection*{You'll want to know: \ja{なげる}}
\textbf{\ja{なげる}} — \emph{nageru} — \textquotedblleft{}to throw\textquotedblright{}.

\begin{quote}
\textbf{\ja{こうえんで} \ja{いぬに} \ja{ぼうを} \ja{なげました。}} — \emph{kōen de inu ni bō o nagemashita} — \textquotedblleft{}In the park I threw a stick for the dog.\textquotedblright{}
\end{quote}

Polite \textbf{\ja{なげます}}, \emph{nagemasu}; an ichidan verb.

\begin{cousinweb}
\textbf{\ja{なげる}} (\emph{nageru}) is a native Japanese word (wago), inherited rather than borrowed from Chinese; its kanji was attached later as a kun'yomi (JMdict, KANJIDIC2).
\end{cousinweb}

\begin{grammarlens}[title={What you've built}]
One more word for everyday life at A2.
\end{grammarlens}

\subsection*{Your turn}
\begin{itemize}
\raggedright
  \item \emph{Say it:} \emph{nageru}
  \item \emph{Say it:} \emph{nageru}, once more
  \item \emph{Say it:} \emph{kōen de inu ni bō o nagemashita} — In the park I threw a stick for the dog.
  \item \emph{Recall:} say the Japanese for a player, then the Japanese for cheering, then say \emph{nageru} again
\end{itemize}

\begin{tcolorbox}[breakable,colback=teal!4,colframe=teal!35!black,title={Before you move on}]
What does \ja{なげる} mean? (\textquotedblleft{}To throw\textquotedblright{}.) Say it once more.
\end{tcolorbox}
\section[\texorpdfstring{\ja{ける} (keru)}{ける (keru)}]{\ja{ける} (keru) — to kick}

\label{lesson:JA-C235-keru}



\begin{quote}
Before the new one: say the Japanese for cheering, then the Japanese for to throw.
\end{quote}

\subsection*{You'll want to know: \ja{ける}}
\textbf{\ja{ける}} — \emph{keru} — \textquotedblleft{}to kick\textquotedblright{}.

\begin{quote}
\textbf{\ja{おとうとは} \ja{にわで} \ja{まりを} \ja{けりました。}} — \emph{otōto wa niwa de mari o kerimashita} — \textquotedblleft{}My younger brother kicked a ball in the garden.\textquotedblright{}
\end{quote}

Polite \textbf{\ja{けります}}, \emph{kerimasu}; a godan verb.

\begin{cousinweb}
\textbf{\ja{ける}} (\emph{keru}) is a native Japanese word (wago), inherited rather than borrowed from Chinese; its kanji was attached later as a kun'yomi (JMdict, KANJIDIC2).
\end{cousinweb}

\begin{grammarlens}[title={What you've built}]
One more word for everyday life at A2.
\end{grammarlens}

\subsection*{Your turn}
\begin{itemize}
\raggedright
  \item \emph{Say it:} \emph{keru}
  \item \emph{Say it:} \emph{keru}, once more
  \item \emph{Say it:} \emph{otōto wa niwa de mari o kerimashita} — My younger brother kicked a ball in the garden.
  \item \emph{Recall:} say the Japanese for cheering, then the Japanese for to throw, then say \emph{keru} again
\end{itemize}

\begin{tcolorbox}[breakable,colback=teal!4,colframe=teal!35!black,title={Before you move on}]
What does \ja{ける} mean? (\textquotedblleft{}To kick\textquotedblright{}.) Say it once more.
\end{tcolorbox}
\section[\texorpdfstring{\ja{やきゅう} (yakyū)}{やきゅう (yakyū)}]{\ja{やきゅう} (yakyū) — baseball}

\label{lesson:JA-C235-yakyuu}



\begin{quote}
Before the new one: say the Japanese for to throw, then the Japanese for to kick.
\end{quote}

\subsection*{You'll want to know: \ja{やきゅう}}
\textbf{\ja{やきゅう}} — \emph{yakyū} — \textquotedblleft{}baseball\textquotedblright{}.

\begin{quote}
\textbf{\ja{きのう} \ja{ともだちと} \ja{やきゅうを} \ja{しました。}} — \emph{kinō tomodachi to yakyū o shimashita} — \textquotedblleft{}Yesterday I played baseball with friends.\textquotedblright{}
\end{quote}

\begin{cousinweb}
\textbf{\ja{やきゅう}} (\emph{yakyū}) is Sino-Japanese (kango): two Chinese characters read with their on'yomi, readings borrowed from Chinese. They are \emph{ya}, \textquotedblleft{}field\textquotedblright{} (Middle Chinese \emph{yaeX}) and \emph{kyū}, \textquotedblleft{}ball\textquotedblright{} (\emph{gjuw}) (KANJIDIC2; Middle Chinese in Baxter's transcription).
\end{cousinweb}

\begin{grammarlens}[title={What you've built}]
One more word for everyday life at A2.
\end{grammarlens}

\subsection*{Your turn}
\begin{itemize}
\raggedright
  \item \emph{Say it:} \emph{yakyū}
  \item \emph{Say it:} \emph{yakyū}, once more
  \item \emph{Say it:} \emph{kinō tomodachi to yakyū o shimashita} — Yesterday I played baseball with friends.
  \item \emph{Recall:} say the Japanese for to throw, then the Japanese for to kick, then say \emph{yakyū} again
\end{itemize}

\begin{tcolorbox}[breakable,colback=teal!4,colframe=teal!35!black,title={Before you move on}]
What does \ja{やきゅう} mean? (\textquotedblleft{}Baseball\textquotedblright{}.) Say it once more.
\end{tcolorbox}
\section[\texorpdfstring{\ja{じゅうどう} (jūdō)}{じゅうどう (jūdō)}]{\ja{じゅうどう} (jūdō) — judo}

\label{lesson:JA-C235-juudou}



\begin{quote}
Before the new one: say the Japanese for to kick, then the Japanese for baseball.
\end{quote}

\subsection*{You'll want to know: \ja{じゅうどう}}
\textbf{\ja{じゅうどう}} — \emph{jūdō} — \textquotedblleft{}judo\textquotedblright{}.

\begin{quote}
\textbf{\ja{あには} \ja{きょねん} \ja{じゅうどうを} \ja{ならいました。}} — \emph{ani wa kyonen jūdō o naraimashita} — \textquotedblleft{}My elder brother took judo lessons last year.\textquotedblright{}
\end{quote}

\begin{cousinweb}
\textbf{\ja{じゅうどう}} (\emph{jūdō}) is Sino-Japanese (kango): two Chinese characters read with their on'yomi, readings borrowed from Chinese. They are \emph{jū}, \textquotedblleft{}soft, gentle\textquotedblright{} (Middle Chinese \emph{nyuw}) and \emph{dō}, \textquotedblleft{}way\textquotedblright{} (\emph{dawX}) (KANJIDIC2; Middle Chinese in Baxter's transcription).
\end{cousinweb}

\begin{grammarlens}[title={What you've built}]
One more word for everyday life at A2.
\end{grammarlens}

\subsection*{Your turn}
\begin{itemize}
\raggedright
  \item \emph{Say it:} \emph{jūdō}
  \item \emph{Say it:} \emph{jūdō}, once more
  \item \emph{Say it:} \emph{ani wa kyonen jūdō o naraimashita} — My elder brother took judo lessons last year.
  \item \emph{Recall:} say the Japanese for to kick, then the Japanese for baseball, then say \emph{jūdō} again
\end{itemize}

\begin{tcolorbox}[breakable,colback=teal!4,colframe=teal!35!black,title={Before you move on}]
What does \ja{じゅうどう} mean? (\textquotedblleft{}Judo\textquotedblright{}.) Say it once more.
\end{tcolorbox}
\section[\texorpdfstring{\ja{ゆうしょう} (yūshō)}{ゆうしょう (yūshō)}]{\ja{ゆうしょう} (yūshō) — a championship win; overall victory}

\label{lesson:JA-C235-yuushou}



\begin{quote}
Before the new one: say the Japanese for baseball, then the Japanese for judo.
\end{quote}

\subsection*{You'll want to know: \ja{ゆうしょう}}
\textbf{\ja{ゆうしょう}} — \emph{yūshō} — \textquotedblleft{}a championship win; overall victory\textquotedblright{}.

\begin{quote}
\textbf{\ja{せんげつ} \ja{あにの} \ja{がっこうが} \ja{ゆうしょうしました。}} — \emph{sengetsu ani no gakkō ga yūshō shimashita} — \textquotedblleft{}Last month my elder brother's school won the championship.\textquotedblright{}
\end{quote}

\begin{cousinweb}
\textbf{\ja{ゆうしょう}} (\emph{yūshō}) is Sino-Japanese (kango): two Chinese characters read with their on'yomi, readings borrowed from Chinese. They are \emph{yū}, \textquotedblleft{}excel\textquotedblright{} (Middle Chinese \emph{'juw}) and \emph{shō}, \textquotedblleft{}win\textquotedblright{} (\emph{syingH}) (KANJIDIC2; Middle Chinese in Baxter's transcription).
\end{cousinweb}

\begin{grammarlens}[title={What you've built}]
One more word for everyday life at A2.
\end{grammarlens}

\subsection*{Your turn}
\begin{itemize}
\raggedright
  \item \emph{Say it:} \emph{yūshō}
  \item \emph{Say it:} \emph{yūshō}, once more
  \item \emph{Say it:} \emph{sengetsu ani no gakkō ga yūshō shimashita} — Last month my elder brother's school won the championship.
  \item \emph{Recall:} say the Japanese for baseball, then the Japanese for judo, then say \emph{yūshō} again
\end{itemize}

\begin{tcolorbox}[breakable,colback=teal!4,colframe=teal!35!black,title={Before you move on}]
What does \ja{ゆうしょう} mean? (\textquotedblleft{}A championship win; overall victory\textquotedblright{}.) Say it once more.
\end{tcolorbox}
\section[(dialogue)]{Words from chapters 230–235 — a review}

\label{lesson:JA-R235-review-230-235}



\begin{quote}
Say the Japanese for a company president and a boss.
\end{quote}

\subsection*{Your turn}
Say these aloud:

\begin{itemize}
\raggedright
  \item a company president, a boss, a senior, a junior, to do one's best
  \item hello on the phone, a number, contact, a response, to pass on
  \item to discuss, to put together, to entrust, to refuse, to inform
  \item to play, to swim, swimming, to run, exercise
  \item a match, to win, to be beaten, a player, cheering
  \item to throw, to kick, baseball, judo, a championship win
\end{itemize}

\begin{tcolorbox}[breakable,colback=teal!4,colframe=teal!35!black,title={Before you move on}]
A company president? (\textbf{\ja{しゃちょう}}, \emph{shachō}.) A boss? (\textbf{\ja{じょうし}}, \emph{jōshi}.) A senior? (\textbf{\ja{せんぱい}}, \emph{senpai}.) A junior? (\textbf{\ja{こうはい}}, \emph{kōhai}.) To do one's best? (\textbf{\ja{がんばる}}, \emph{ganbaru}.) Hello on the phone? (\textbf{\ja{もしもし}}, \emph{moshimoshi}.) A number? (\textbf{\ja{ばんごう}}, \emph{bangō}.) Contact? (\textbf{\ja{れんらく}}, \emph{renraku}.) A response? (\textbf{\ja{へんじ}}, \emph{henji}.) To pass on? (\textbf{\ja{つたえる}}, \emph{tsutaeru}.) To discuss? (\textbf{\ja{はなしあう}}, \emph{hanashiau}.) To put together? (\textbf{\ja{まとめる}}, \emph{matomeru}.) To entrust? (\textbf{\ja{まかせる}}, \emph{makaseru}.) To refuse? (\textbf{\ja{ことわる}}, \emph{kotowaru}.) To inform? (\textbf{\ja{しらせる}}, \emph{shiraseru}.) To play? (\textbf{\ja{あそぶ}}, \emph{asobu}.) To swim? (\textbf{\ja{およぐ}}, \emph{oyogu}.) Swimming? (\textbf{\ja{すいえい}}, \emph{suiei}.) To run? (\textbf{\ja{はしる}}, \emph{hashiru}.) Exercise? (\textbf{\ja{うんどう}}, \emph{undō}.) A match? (\textbf{\ja{しあい}}, \emph{shiai}.) To win? (\textbf{\ja{かつ}}, \emph{katsu}.) To be beaten? (\textbf{\ja{まける}}, \emph{makeru}.) A player? (\textbf{\ja{せんしゅ}}, \emph{senshu}.) Cheering? (\textbf{\ja{おうえん}}, \emph{ōen}.) To throw? (\textbf{\ja{なげる}}, \emph{nageru}.) To kick? (\textbf{\ja{ける}}, \emph{keru}.) Baseball? (\textbf{\ja{やきゅう}}, \emph{yakyū}.) Judo? (\textbf{\ja{じゅうどう}}, \emph{jūdō}.) A championship win? (\textbf{\ja{ゆうしょう}}, \emph{yūshō}.)
\end{tcolorbox}
